\section{Testing error-correcting codes}\label{sec:error}

In \Cref{sec:answer-reduction} below,
rather than the prover sending the verifier their entire ``large" answer~$a$,
they will instead encode it into $\codee(a)$ using an error correcting code
and allow the verifier to query individual bits of the encoding.
(The fact that the verifier is allowed to query bits of the encoding rather than the original string
stems from the PCPP technology we use. See \Cref{sec:composing-with-an-error-correcting-code} for more details.)
In this section, we develop the tests which verify that provers are performing this task honestly,
so that when we query a subset of the bits~$I$,
they respond based on the bits of a codeword which was sampled independently of~$I$.
We will develop such a test for the low-degree code (\Cref{sec:testing-low-deg}).

Our proofs are entirely standard: we start with the known property tester for this code (i.e..\ \Cref{thm:anand-thomas-classical-low-degree}),
which allows us to query the prover's codeword at a uniformly random location.
Then we use the local decodability properties of this code to allow us to query arbitrary subsets of coordinates.
We begin by stating a slightly nonstandard definition of error-correcting codes relevant to our application.

\begin{definition}[Error-correcting codes]\label{def:error-correcting-code}
Let $m$ and $q$ be integers,
and let $\eta \in [0, 1]$.
An \emph{$(n, m, q, \eta)$-error-correcting code} $\codec = (\codee,\coded,\codes)$ is defined as follows.
\begin{itemize}
\item[$\circ$] $\codes$ is a subset of $\F_q^m$ such that for each $x \neq y \in \codes$, $x$ and $y$ have normalized Hamming agreement at most~$\eta$
		(i.e.\ the probability, over a uniformly random $\bi \in [m]$, that $x_{\bi} = y_{\bi}$ is at most $\eta$).
\item[$\circ$] $\codee:\{0, 1\}^n \rightarrow \codes \subseteq \F_q^m$ is the \emph{encoding map}.
\item[$\circ$] $\coded: \F_q^m \rightarrow \{0, 1\}^n \cup \{\bot\}$ is the \emph{decoding map}.
		For each $x \in \{0, 1\}^n$, $\coded(\codee(x)) = x$.
		In addition, for every~$w$ not in the range of~$\codee$, $\coded(w) = \bot$.
\end{itemize}
\end{definition}

\begin{remark}
The purpose of the subset $\codes$ is this:
in this section, we are designing games which test that a prover responds according to an error-correcting code.
This means that the prover should respond based on the encoding $\codee(x)$ of some string~$x \in \{0, 1\}^n$.
However, the games we design may only be able to test if the prover responds based on a string in~$\codes$,
which contains the encodings $\codee(x)$ but may include other strings as well.
This definition ensures that these other strings are still far from each other in Hamming distance.
\end{remark}

The next definition defines a subset tester.

\begin{definition}\label{def:testing-a-code}
Let $\codec = (\codee, \coded, \codes)$ be an $(n, m, q, \eta)$-error-correcting code.
Let~$k$ be an integer.
Given a game~$\game(\cdot)$  whose inputs are from the set of subsets of~$[m]$ of size~$k$
and a probability distribution~$\calD$ over this set,
we write $\game(\calD)$ for the game in which we first sample $\bI \sim \calD$ and then play~$\game(\bI)$.
Then $\game$ is a \emph{$k$-subset tester with robustness $\delta(\eps)$} for $\codec$ if it satisfies the following two properties.
\begin{itemize}
\item[$\circ$] \textbf{Completeness:}
Let $(\psi, M)$ be an EPR strategy in which $\{M_w\}$ is a measurement with outcomes in~$\{0, 1\}^n$.
Consider the partial strategy $(\psi, G)$ in which
\begin{equation*}
	G^I_{a_1, \ldots, a_k} := M_{[\codee(w)|_I = a_1, \ldots, a_k]}.
\end{equation*}
Then this can be extended to a (full) real commuting EPR strategy which, for each~$I$, passes $\game(I)$ with probability~$1$.
\item[$\circ$] \textbf{Soundness:} For any distribution~$\calD$, consider a strategy $(\psi, M)$ which passes $\game(\calD)$ with probability $1-\eps$.
Then there exists a measurement $\{G_w\}_{w}$ with outcomes~$w$ in~$\codes$ such that
\begin{equation*}
	M^I_{a_1, \ldots, a_k} \otimes I_{\reg{Bob}} \consistency_{\delta(\eps)} I_{\reg{Alice}} \otimes G_{[w|_I = a_1, \ldots, a_k]}.
\end{equation*}
\end{itemize}
\end{definition}

\subsection{Testing the low-degree code}\label{sec:testing-low-deg}

In this section, we show how to test the low-degree code.
This is essentially an exercise in generalizing \Cref{thm:anand-thomas-classical-low-degree} to arbitrary subsets.
We begin with some notation.

\begin{notation}
We write $\mathcal{F}_{q, k}^{m}$ for the family
$\displaystyle
\mathcal{F}^m_{q, k} = \{F \subseteq \F_q^m \mid |F| \leq k\}
$.
\end{notation}

Now we define the low-degree code tester.

\begin{definition}
Let $m$, $q$, and~$d$ be integers.  Let~$k$ be an integer, and let~$F$ be an element of~$\mathcal{F}^m_{q, k}$.
Then $\game_{\mathrm{LDsubset}}(m, q, d, F)$ is the game defined in \Cref{fig:classical-subset}.
\end{definition}

{
\floatstyle{boxed} 
\restylefloat{figure}
\begin{figure}
With probability~$\tfrac{1}{2}$ each, perform one of the following two tests.
\begin{enumerate}
	\item \textbf{Low-degree:} Perform $\game_{\mathrm{Surface}}(m, d, q,2)$. \label{item:low-deg-point}
	\item \textbf{Cross-check:}
			Flip an unbiased coin $\bb \sim \{0, 1\}$.
			Let~$\bs$ be a uniformly random subspace of dimension $k+1$ containing the points in~$F$.
			 With probability $\tfrac{1}{2}$ each:
			\begin{enumerate}
			\item Let $\bw$ be a uniformly random point in~$\bs$. Distribute the questions as follows:\label{item:bootstrapping-to-subspaces}
				\begin{itemize}
				\item[$\circ$] Player~$\bb$: give $\bw$; receive a value $\by \in \F_q$.
				\item[$\circ$] Player~$\overline{\bb}$: give $\bs$; receive a degree-$d$ polynomial $\bg:\bs \rightarrow \F_q$.
				\end{itemize}
				Accept if $\bg(\bw) = \by$.
			\item Distribute the questions as follows: \label{item:boosting-to-subsets}
				\begin{itemize}
				\item[$\circ$] Player~$\bb$: give $\bs$; receive a degree-$d$ polynomial $\bg:\bs \rightarrow \F_q$.
				\item[$\circ$] Player~$\overline{\bb}$: give $F$; receive a function $\boldf:\bF \rightarrow \F_q$.
				\end{itemize}
				Accept if $\bg|_{F} = \boldf$.
			\end{enumerate}
	\end{enumerate}
	\caption{The game $\game_{\mathrm{LDsubset}}(m, q, d, F)$. \label{fig:classical-subset}}
\end{figure}
}

The performance of the low-degree subset game is given by the following theorem.

\begin{theorem}\label{thm:low-degree-code-tester}
Consider low-degree parameters $\params = (n, q, h, H, m, \calS, \pi)$. Set $d = m(h-1)$.
Set $m' = q^m$.  We will identify strings in $\F_q^{m'}$ with functions $g:\F_q^m \rightarrow \F_q$.
Given $a \in \{0, 1\}^n$, define $\codee(a) = g_a$ and $\coded(g_a) = a$.
For all other $g:\F_q^m \rightarrow \F_q$ (i.e.\ those which are not the low-degree encoding of a string~$a$),
define $\coded(g) = \bot$.
Finally, define $\codes$ to be the set of degree~$d$ polynomials~$g:\F_q^m \rightarrow\F_q$.
Then $\codec = (\codee, \coded, \codes)$ is an $(n, m', q, d/q)$-error-correcting code.

Furthermore, there exists a constant $c > 0$ and a function $\delta(\eps) = \poly(\eps, dm/q^c)$ such that the following holds.
Let~$k$ be an integer.
Then $\game_{\mathrm{LDsubset}} := \game_{\mathrm{LDsubset}}(m, q, d, \cdot)$ is a $k$-subset tester for $\mathrm{LDCode}$ with robustness $\delta(\eps)$.

Finally,
\begin{equation*}
\qtime{\game_{\mathrm{LDsubset}}} = \poly(m, k,\log q),
\quad
\atime{\game_{\mathrm{LDsubset}}} = \poly(m, d^k, \log q),
\end{equation*}
\begin{equation*}
\qlength{\game_{\mathrm{LDsubset}}} = O(k m \log q),
\quad
\alength{\game_{\mathrm{LDsubset}}} = O(d^k \log (q)).
\end{equation*}
\end{theorem}

Before proving this, we need the following proposition.

\begin{proposition}\label{prop:almost-uniform}
Let $F \subseteq \F_q^m$ be of size at most~$k$. Consider the distribution $\calD_{\mathrm{twostep}}$ on points~$x\in\F_q^m$ generated by the following two-step process:
(i) let $\bs$ be a uniformly random subspace of size~$k+1$ containing~$F$, and (ii) draw $\bx$ uniformly at random from~$\bs$.
Let $\calD_{\mathrm{unif}}$ be the uniform distribution on~$\F_q^m.$
Then $d_\mathrm{TV}(\calD_{\mathrm{twostep}}, \calD_{\mathrm{unif}}) \leq 1/q$.
\end{proposition}
\begin{proof}
Let $x_1, \ldots, x_\ell$ be a maximal set of linearly independent elements from~$F$.
A uniformly random subspace of size~$k+1$ containing~$F$ can be generated as follows:
first, pick a uniformly random nonzero vector $\by_{\ell+1}$ linearly independent from~$F$,
then pick a uniformly random nonzero vector $\by_{\ell+2}$ linearly independent from~$F\cup\{\by_{\ell+1}\}$,
and so forth. Set $\bs = \mathrm{span}\{x_1, \ldots, x_\ell, \by_{\ell+1}, \ldots, \by_{k+1}\}$.
A uniformly random point in~$\bs$ will be of the form 
\begin{equation*}
\bx = \bt_1 x_1 + \cdots + \bt_\ell x_\ell + \bt_{\ell+1} \by_{\ell+1} + \cdots + \bt_{k+1} \by_{k+1},
\end{equation*}
where each $\bt_{i}$ is a uniformly random element in $\F_q$.
Because all the $\by_i$'s are linearly independent,
the linear combination $\bt_{\ell+1} \by_{\ell+1} + \cdots + \bt_{k+1} \by_{k+1}$
is zero only when $\bt_{\ell+1} = \cdots = \bt_{k+1} = 0$.
Otherwise, this linear combination is distributed as a uniformly random nonzero vector linearly independent from~$F$.
Thus, with probability $(q^{k+1-\ell})^{-1}$, $\bx$ is distributed as a uniformly random vector in the span of~$F$,
and otherwise it is distributed as a uniformly random vector outside the span of~$F$.
Given that the span of~$F$ has $q^{\ell}$ points, the total variation distance is
\begin{equation*}
\frac{1}{2} \left| \frac{q^{\ell}}{q^m} - \frac{1}{q^{k+1-\ell}}\right|
+ \frac{1}{2} \left| \frac{q^m- q^{\ell}}{q^m} - \left(1- \frac{1}{q^{k+1-\ell}}\right)\right|
= q^{\ell}\left(\frac{1}{q^{k+1}} - \frac{1}{q^m}\right)
\leq \frac{1}{q}.\qedhere
\end{equation*}
\end{proof}

Now we prove \Cref{thm:low-degree-code-tester}.

\begin{proof}[Proof of \Cref{thm:low-degree-code-tester}]
The fact that $\codec$ is an $(n, m', q, d/q)$-error-correcting code follows from Schwartz-Zippel (\Cref{lem:schwartz-zippel}).

\paragraph{Completeness.}
Let $(\psi, M)$ be an EPR strategy in which $\{M_w\}$ is a measurement with outcomes in~$\{0, 1\}^n$.
Consider the strategy $(\psi, G)$ in which for any subset of points~$F = \{y_1, \ldots, y_\ell\}$,
\begin{equation*}
	G^I_{a_1, \ldots, a_\ell} := M_{[g_x(y_1), \ldots, g_x(y_\ell) = a_1, \ldots, a_\ell]}.
\end{equation*}
(This covers the case of points ($\ell = 1$) and subsets~$F$ ($\ell = k$).)
In addition, for any subspace~$s$,
\begin{equation*}
	G^s_{f} := M_{[g_x|_s = f]}.
\end{equation*}
(This covers the case of the $2$-dimensional subspaces used in $\game_{\mathrm{Point}}$
and the $(k+1)$-dimensional subspaces used for the local decoding.)
By construction, $(\psi, G)$ is an EPR strategy,
and it is easy to see that it is a \emph{commuting} one as well.

We claim that $(\psi, G)$ passes $\game_{\mathrm{LDsubset}}(m, q, d, \calD)$ with probability~$1$.
We begin with the low-degree test.
By \Cref{fact:heh-heh-heh-gonna-make-anand-prove-this-so-i-can-take-the-day-off},
$M_w \otimes I_{\reg{Bob}} \consistency_0 I_{\reg{Alice}} \otimes M_w$.
Then by \Cref{fact:specialize-the-simeq},
\begin{equation*}
M_{[g_x(w) = b]} \otimes I_{\reg{Bob}}
\consistency_{0} I_{\reg{Alice}} \otimes M_{[g_x(w) = b]}.
\end{equation*}
This implies passing the low-degree test with probability~$1$, because
\begin{equation*}
G^s_{[f(w) = b]} \otimes I_{\reg{Bob}}
= M_{[g_x(w) = b]} \otimes I_{\reg{Bob}}
\consistency_{0} I_{\reg{Alice}} \otimes M_{[g_x(w) = b]} = I_{\reg{Alice}} \otimes G^w_b.
\end{equation*}
A similar argument shows the other tests pass with probability~$1$ as well.

\paragraph{Soundness.}

Let~$\calD$ be a distribution,
and let $(\psi, M)$ be a strategy which passes $\game_{\mathrm{LDsubset}}(m, q, d, \calD)$ with probability $1-\eps$.
The outline of the proof is as follows:
first we will use the low degree test in \Cref{item:low-deg-point} to ensure the test correctly answers low-degree point queries.
\Cref{item:bootstrapping-to-subspaces} will then bootstrap this to subspaces,
and \Cref{item:boosting-to-subsets} will further bootstrap this to subsets, proving the theorem.

\paragraph{Using the low-degree test.}
Passing the test with probability $1-\epsilon$ means passing the low-degree test with probability at least $1-2\epsilon$.
By \Cref{thm:anand-thomas-classical-low-degree}, this means that there exists a POVM measurement $G \in \polymeas{m}{d}{q}$ such that
\begin{equation}\label{eq:just-applied-low-degree}
M^w_b \otimes I_{\reg{Bob}} \consistency_{\delta(\eps)} I_{\reg{Alice}} \otimes G_{[g(w) = b]},
\qquad G_g \otimes I_{\reg{Bob}} \consistency_{\delta(\eps)} I_{\reg{Alice}} \otimes G_g,
\end{equation}
where the first is on the uniform distribution over $\F_q^m$.

\paragraph{Bootstrapping to subspaces.}
Define $\calD_{\mathrm{twostep}}$ to be the two-step sampling process $(\bF, \bs, \bw)$ as in \Cref{item:bootstrapping-to-subspaces}.
By \Cref{prop:almost-uniform}, the marginal distribution on~$\bw$ has total variation distance at most~$1/q$ with $\calD_{\mathrm{unif}}$.
As a result, we can apply \Cref{fact:swap-dists} to \Cref{eq:just-applied-low-degree}, yielding
\begin{equation}\label{eq:gonna-triangle-the-shit-outta-this}
M^w_b \otimes I_{\reg{Bob}} \consistency_{\delta(\eps)} I_{\reg{Alice}} \otimes G_{[g(w) = b]}
\end{equation}
on distribution $\calD_{\mathrm{twostep}}$.
Similarly, by \Cref{fact:specialize-the-simeq},
\begin{equation}\label{eq:incoming-triangle}
G_{[g(w) = b]} \otimes I_{\reg{Bob}} \consistency_{\delta(\eps)} I_{\reg{Alice}} \otimes G_{[g(w) = b]}
\end{equation}
on distribution $\calD_{\mathrm{twostep}}$.


Next, because the strategy passes the test in \Cref{item:bootstrapping-to-subspaces} with probability at least $1-4\epsilon$,
\begin{equation}\label{eq:watch-out-for-the-triangle}
M^{w}_y \otimes I_{\reg{Bob}} \consistency_{\epsilon} I_{\reg{Alice}} \otimes M^s_{[g(w) = y]}.
\end{equation}
on distribution $\calD_{\mathrm{twostep}}$.
Combining \Cref{eq:gonna-triangle-the-shit-outta-this,eq:incoming-triangle,eq:watch-out-for-the-triangle} with our second triangle inequality (\Cref{fact:triangle-like}),
\begin{equation*}
M^s_{[g(w) = y]} \otimes I_{\reg{Bob}}
\consistency_{\delta(\eps)} I_{\reg{Alice}} \otimes G_{[g(w) = b]}.
\end{equation*}
By~\Cref{prop:same-on-point-same-on-subspace}, we conclude that
\begin{equation}\label{eq:finished-with-subspaces}
M^s_{f} \otimes I_{\reg{Bob}}
\consistency_{\delta(\eps)} I_{\reg{Alice}} \otimes G_{[g|_s = f]}.
\end{equation}

\paragraph{Concluding with subsets.}
The strategy passes the test in \Cref{item:boosting-to-subsets} with probability at least $1-4\epsilon$. As a result,
\begin{equation*}
M^F_f \otimes I_{\reg{Bob}}
\consistency_{\epsilon}
I_{\reg{Alice}} \otimes M^s_{[g|_F = f]}.
\end{equation*}
Applying \Cref{fact:specialize-the-simeq} to \Cref{eq:finished-with-subspaces},
\begin{equation*}
M^s_{[g|_F=f]} \otimes I_{\reg{Bob}}
\consistency_{\delta(\eps)} I_{\reg{Alice}} \otimes G_{[h|_F = h]}.
\end{equation*}
Similarly, applying \Cref{fact:specialize-the-simeq} to \Cref{eq:just-applied-low-degree},
\begin{equation*}
G_{[h|_F = f]} \otimes I_{\reg{Bob}} \consistency_{\delta(\eps)} I_{\reg{Alice}} \otimes G_{[h|_F=f]}
\end{equation*}
Applying the triangle inequality (\Cref{fact:triangle-like}) to these three equations,
we get
\begin{equation*}
M^F_f \otimes I_{\reg{Bob}}
\consistency_{\delta(\epsilon)}
I_{\reg{Alice}} \otimes G_{[g|_{F} = f]}
\end{equation*}
with respect to distribution $\calD_{\mathrm{twostep}}$,
and therefore, by \Cref{fact:close-on-marginal}, with respect to $\calD$.
\end{proof}

\subsection{Efficiently decodable codes}

Our application requires error-correcting codes with two further properties.
The first property is that the decoding map $\coded(\cdot)$ be efficiently computable.
(The encoding map, on the other hand, is allowed arbitrary complexity.
This is because we will leave the task of computing the encoding maps to the provers.) 
The second, more technical property is
we require that the code \emph{embed} the codeword, in the following sense:
the encoding $\codee(x)$ of a string~$x$ should actually \emph{contain} the string~$x$,
and the function for where to find each bit of~$x$ in $\codee(x)$ should be efficiently computable.

\begin{definition}[Efficiently-decodable error-correcting codes]\label{def:efficient-codes}
Let $m, q:\Z^+\rightarrow \Z^+$,
and let $\eta:\Z^+\rightarrow [0, 1]$.
Let $t_{\mathrm{Dec}}, t_{\mathrm{Emb}}:\Z^+ \rightarrow \Z^+$.
We say that $\codec_n=(\codee_n, \coded_n, \codes_n)$ is an \emph{$(n, m, q, \eta, t_{\mathrm{Dec}}, t_{\mathrm{Emb}})$-efficient code family}
if the following three conditions are true.
\begin{itemize}
\item[$\circ$] For each~$n$, $(\codee_n, \coded_n, \codes_n)$ is an $(n, m(n), q(n), \eta(n))$-error-correcting code.
\item[$\circ$] There exists an algorithm $\mathrm{Alg}_{\coded}$ which, on input $(n, w)$, outputs $\coded_n(w)$.
			Furthermore, $\mathrm{Alg}_{\coded}$ runs in time $t_{\mathrm{Dec}}(n)$.
\item[$\circ$] There exists an embedding $\mu_n:[n] \rightarrow [m(n)]$ such that for each $i \in [n]$, $x_i = (\codee_n(x))_{\mu_n(i)}$.
			Furthermore, there is an algorithm $\mathrm{Alg}_{\mathrm{Emb}}$ which, on input $(n, i)$, computes $\mu_n(i)$ in time
			 $t_{\mathrm{Emb}}(n)$.
\end{itemize}
\end{definition}

Now, we show that the low-degree code is efficiently-decodable.
The decoding algorithm follows a simple strategy:
assuming that the input is a proper encoding of a message,
they can directly read off the message from the input.
Then they compute the encoding of the purported message and check that it equals the input.

\begin{fact}\label{fact:low-degree-code-is-efficient}
There is a $(n, m', q, \eta, t_{\mathrm{Dec}}, t_{\mathrm{Emb}})$-error-correcting code $\codec$ with parameters set as follows:
\begin{equation*}
m'(n) = \poly(n), \quad q(n) = \mathrm{polylog}(n), \quad \eta(n) = \frac{1}{\mathrm{polylog}(n)},
\end{equation*}
\begin{equation*}
t_{\mathrm{Dec}}(n) = \poly(n), \qquad t_{\mathrm{Emb}}(n) = \mathrm{polylog}(n).
\end{equation*}
In addition, $\codec$ has a $k$-subset test~$\game$ with robustness $\delta(\eps) = \poly(\eps, 1/\log(n))$ such that
\begin{equation*}
\qtime{\game} = \poly(\log n, k),
\quad
\atime{\game} = \poly(\log(n)^k),
\end{equation*}
\begin{equation*}
\qlength{\game} = O(k \log n),
\quad
\alength{\game} = O(\log(n)^{2k}).
\end{equation*}
\end{fact}
\begin{proof}
We instantiate the canonical low-degree encoding from \Cref{def:canonical-low-degree} with the ``rule of thumb" parameters from \Cref{eq:parameters}:
\begin{equation*}
h(n) = \Theta(\log(n)),
\qquad
m(n) = \Theta\left(\frac{\log(n)}{\log\log(n)}\right),
\qquad
q(n) = \mathrm{polylog}(n).
\end{equation*}
If we set $d(n) = m(n)\cdot (h(n) - 1)$, then this is a code with distance $\eta(n) = 1-d(n)/q(n) = 1-1/\mathrm{polylog}(n)$.
In addition, it has length $m'(n) = q(n)^{m(n)} = \poly(n)$.
Finally, the canonical low-degree encoding gives us the embedding $\mu_{\mathrm{Emb}} := \sigma_{m, t_1, t_2}$.
By \Cref{prop:canonical-time}, it takes time $t_{\mathrm{Emb}}(n) = \mathrm{poly log}(n)$ to compute.

Now we design the decoding algorithm $\mathrm{Alg}_{\coded}$.
On input $(n, w)$, it rejects if~$w$ is not length~$m'$.
Otherwise, it interprets~$w$ as a function $f:\F_q^m\rightarrow \F_q$.
It queries~$g$ on the points $\pi(1), \ldots, \pi(n)$.
Let $a \in \{0, 1\}^n$ be the received answers.
If~$g$ is a codeword, it equals the low-degree function~$g_a$.
So the algorithm simply iterates over all~$x \in \F_q^m$ and checks that $f(x) = g_a(x)$.
By \Cref{prop:canonical-time}, computing $g_a(x)$ can be done in time $\poly(n)$,
and so this takes time $t_{\mathrm{Dec}}(n) = \poly(n)$ in total.

Finally, the performance of the subset tester follows from \Cref{thm:low-degree-code-tester} with our setting of parameters.
\end{proof}


%%% Local Variables:
%%% mode: latex
%%% TeX-master: "main"
%%% End:
%%%---------- close: error-correcting-code-test

%%%%%%%%%%% jump to answer-reduction
%%%---------- open: answer-reduction
%!TEX root = main.tex

